%!TEX program = xelatex
%!TEX encoding = ISO-8859-1
\documentclass[12pt]{article}

\usepackage{amssymb}
\usepackage{amsmath,amsfonts,amsthm,amstext,mathabx}
\usepackage[brazil]{babel}
%\usepackage[latin1]{inputenc}
\usepackage{graphicx}
\graphicspath{{/home/jfreitas/Dropbox/imagens-latex/}{/ArquivosUbuntu/Dropbox/imagens-latex/}{D:/Dropbox/imagens-latex/}}
\usepackage{enumitem}
\usepackage{multicol}
\usepackage[all]{xy}

\setlength{\topmargin}{-1.0in}
\setlength{\oddsidemargin}{0in}
\setlength{\textheight}{10.1in}
\setlength{\textwidth}{6.5in}
\setlength{\baselineskip}{12mm}

\newcounter{exercicios}
\renewcommand{\theexercicios}{\arabic{exercicios}}

\newenvironment{questao}[1]{%
\refstepcounter{exercicios}%
\label{#1}%
\noindent\textbf{Exerc{\'\i}cio {\theexercicios}:}%
}%
{}%

% \newcounter{exercicios}
% \setcounter{exercicios}{0}
% \newcommand{\questao}{
% \addtocounter{exercicios}{1}
% \noindent{\bf Exerc{\'\i}cio \arabic{exercicios}: }}

\newcommand{\equi}{\Leftrightarrow}
\newcommand{\bic}{\leftrightarrow}
\newcommand{\cond}{\rightarrow}
\newcommand{\impl}{\Rightarrow}
\newcommand{\nao}{\sim}
\newcommand{\sub}{\subseteq}
\newcommand{\e}{\ \wedge\ }
\newcommand{\ou}{\ \vee\ }
\newcommand{\vaz}{\emptyset}
\newcommand{\nsub}{\nsubseteq}
\renewcommand{\sin}{{\rm sen\,}}

\newcommand{\n}{\mathbb{N}}
\newcommand{\z}{\mathbb{Z}}
\newcommand{\q}{\mathbb{Q}}
\newcommand{\im}{{\rm Im\,}}
\newcommand{\aut}{{\rm Aut\,}}
\newcommand{\cp}[1]{\mathbb{#1}}
\newcommand{\real}{\mathbb{R}}
\newcommand{\complex}{\mathbb{C}}
\newcommand{\lap}[1]{\mathcal{L}\left\{#1\right\}}
\newcommand{\lapi}[1]{\mathcal{L}^{-1}\left\{#1\right\}}
\newcommand{\se}[1]{\displaystyle\sum_{n = 1}^\infty{#1}}
\newcommand{\dlim}[2]{\displaystyle\lim_{#1\rightarrow #2}}
\newcommand{\slim}{\displaystyle\lim_{n \rightarrow \infty}}
\newcommand{\seq}[1]{\{{#1_n\}}}
\newcommand{\seg}[1]{\displaystyle\sum_{n = 1}^\infty{#1_n}}
\newcommand{\sei}[2]{\displaystyle\sum_{#1}^\infty{#2}}
\newcommand{\sepc}[3]{\displaystyle\sum_{#1}^\infty{#2(x - #3)^n}}
\newcommand{\imp}[3]{\displaystyle\int_{#1}^{+\infty}{#3}{d #2}}
\newcommand{\dint}[4]{\displaystyle\int_{#1}^{#2}{#4}{d#3}}
\newcommand{\inti}[2]{\displaystyle\int{#1}{d#2}}
\newcommand{\norma}[1]{\left\lVert#1\right\rVert}
\newcommand{\flim}[1]{\displaystyle\lim_{#1\rightarrow \infty}}
\renewcommand{\sin}{{\rm sen\,}}
\renewcommand{\tan}{{\rm tg\,}}
\renewcommand{\csc}{{\rm cossec\,}}
\renewcommand{\cot}{{\rm cotg\,}}
\renewcommand{\sinh}{{\rm senh\,}}
\newcommand\T{\rule{0pt}{2.6ex}} 
\newcommand{\vesp}{\vspace{0.2cm}}
\newcommand{\subne}{\subsetneqq}


\newcommand{\compcent}[1]{\vcenter{\hbox{$#1\circ$}}}
\newcommand{\comp}{\mathbin{\mathchoice
{\compcent\scriptstyle}{\compcent\scriptstyle}
{\compcent\scriptscriptstyle}{\compcent\scriptscriptstyle}}}

\begin{document}

\pagestyle{empty}

\begin{figure}[h]
        \begin{minipage}[c]{1.7cm}
        \includegraphics[width=1.7cm]{unb.pdf}
        \end{minipage}%
        \hspace{0pt}
        \begin{minipage}[c]{4in}
          {Universidade de Bras{\'\i}lia} \\
          {Departamento de Matem{\'a}tica}
\end{minipage}
\end{figure}
\vspace{-1cm}\hrule


\begin{center}
{\Large\bf {\'A}lgebra 1 - Turma B -- 2$^{o}$/2017} \\ \vspace{9pt} {\large\bf
  $1^{\underline{o}}$ Teste - Resolu\c{c}\~ao}\\
\vspace{9pt} Prof. Jos{\'e} Ant{\^o}nio O. Freitas
\end{center}
\hrule

\vspace{.6cm}

% \textbf{O Exerc{\'\i}cio 1 foi corrigido para atribui\c{c}\~ao da nota.}

\vspace{.6cm}

\questao Sejam $A$ e $B$ conjuntos. Suponha que $x \in A$ e $A \nsub B$. Então $x \notin B$?

\noindent\textbf{Solu\c{c}\~ao:} Falso. Tome por exemplo $A = \{1,2,3\}$, $B = \{3,4\}$ e $x = 3$. Temos $A \nsub B$, $x \in A$ e no entanto $x \in B$.

\vspace{.5cm}

\questao Mostre que $A \cup ( B \cap C) = (A \cup B) \cap (A \cup C)$.

\noindent\textbf{Solu\c{c}\~ao:} Para mostrar a igualdade desses conjuntos precisamos mostrar que:
\begin{enumerate}
	\item [$1^o$)] $A \cup ( B \cap C) \sub (A \cup B) \cap (A \cup C)$

	Seja $x \in A \cup ( B \cap C)$. Ent\~ao $x \in A$ ou $x \in (B \cap C)$. Se $x \in A$, ent\~ao $x \in A \cup B$ e $x \in (A \cup C)$. Assim $x \in (A \cup B) \cap (A \cup C)$. Com isso $A \cup ( B \cap C) \sub (A \cup B) \cap (A \cup C)$. Agora suponha que $x \in (B \cap C)$. Ent\~ao $x \in B$ e $x \in C$. Neste caso $x \in A \cup B$ e $x \in (A \cup C)$. Logo $x \in (A \cup B) \cap (A \cup C)$ e novamente $A \cup ( B \cap C) \sub (A \cup B) \cap (A \cup C)$. Portanto, independente do caso, sempre temos $A \cup ( B \cap C) \sub (A \cup B) \cap (A \cup C)$.

	\item [$2^o$)] $(A \cup B) \cap (A \cup C) \sub A \cup ( B \cap C)$

	Seja $x \in (A \cup B) \cap (A \cup C)$. Assim $x \in A \cup B$ e $x \in A \cup C$. Daí $x \in A$ ou $x \in B$ e $x \in A$ ou $x \in C$. Se $x \in A$, então $x \in A \cup (B \cap C)$. Logo $(A \cup B) \cap (A \cup C) \sub A \cup ( B \cap C)$. Agora, se $x \notin A$, então $x \in B$ e $x \in C$, isto é, $x \in B \cap C$. Com isso $x \in A \cup (B \cap C)$ e então $(A \cup B) \cap (A \cup C) \sub A \cup ( B \cap C)$. Portanto, independente do caso, $(A \cup B) \cap (A \cup C) \sub A \cup ( B \cap C)$.
\end{enumerate}
Com isso provamos que $A \cup ( B \cap C) = (A \cup B) \cap (A \cup C)$, como quer{\'\i}amos.

\end{document}